% -----------------------------------------------------------------------------
% Class-space remap tables for the PASTIS -> NRW cross-domain inference pipeline
% (companion to pastis_to_germany_inference.py and pastis_to_germany_article.txt).
%
% Paste into a new Overleaf project and compile with pdfLaTeX.
% -----------------------------------------------------------------------------
\documentclass[11pt]{article}
\usepackage[a4paper, margin=2.2cm]{geometry}
\usepackage[table]{xcolor}
\usepackage{booktabs}
\usepackage{array}
\usepackage{caption}

% Common-8 colors, matching COMMON_CMAP in pastis_to_germany_inference.py
\definecolor{cNonCrop}{HTML}{555555}
\definecolor{cWheat}{HTML}{FF7F0E}
\definecolor{cBarley}{HTML}{2CA02C}
\definecolor{cCorn}{HTML}{FFBB78}
\definecolor{cRapeseed}{HTML}{8C564B}
\definecolor{cGrassland}{HTML}{AEC7E8}
\definecolor{cOtherCrop}{HTML}{9467BD}
\definecolor{cVoid}{HTML}{FFFFFF}

% Bordered color swatch; visible even for the white "void" entry.
\newcommand{\swatch}[1]{%
  \fcolorbox{black!70}{#1}{\rule[-0.05em]{0pt}{0.8em}\hspace{1.2em}}%
}

% Tightened column types
\newcolumntype{C}{>{\centering\arraybackslash}m{0.5cm}}
\newcolumntype{L}{>{\raggedright\arraybackslash}m{4.6cm}}

\begin{document}

\section*{Class-space remap tables}

The PASTIS-trained ConvLSTM emits predictions in PASTIS's 20-class label
space; EuroCrops DE\_NRW\_2021 ships labels in the HCAT taxonomy. Both
are mapped into a single 8-class common crop set
(\textit{non\_crop, wheat, barley, corn, rapeseed, grassland, other\_crop, void})
so that cross-domain accuracy can be reported on a shared label set.
Table~\ref{tab:pastis-to-common} gives the prediction-side mapping;
Table~\ref{tab:hcat-to-common} gives the label-side mapping.

% -----------------------------------------------------------------------------
\begin{table}[h]
\centering
\caption{PASTIS 20-class $\to$ common 8-class lookup
(\texttt{PASTIS\_TO\_COMMON}). Class indices match
\texttt{torchgeo.datasets.PASTIS.classes}.}
\label{tab:pastis-to-common}
\small
\setlength{\tabcolsep}{6pt}
\begin{tabular}{C l C C l}
\toprule
\multicolumn{2}{c}{PASTIS class}
& \multicolumn{3}{c}{Common class} \\
\cmidrule(lr){1-2}\cmidrule(lr){3-5}
idx & name & idx & color & name \\
\midrule
0  & background                & 0 & \swatch{cNonCrop}   & non\_crop \\
1  & meadow                    & 5 & \swatch{cGrassland} & grassland \\
2  & soft\_winter\_wheat       & 1 & \swatch{cWheat}     & wheat \\
3  & corn                      & 3 & \swatch{cCorn}      & corn \\
4  & winter\_barley            & 2 & \swatch{cBarley}    & barley \\
5  & winter\_rapeseed          & 4 & \swatch{cRapeseed}  & rapeseed \\
6  & spring\_barley            & 2 & \swatch{cBarley}    & barley \\
7  & sunflower                 & 6 & \swatch{cOtherCrop} & other\_crop \\
8  & grapevine                 & 6 & \swatch{cOtherCrop} & other\_crop \\
9  & beet                      & 6 & \swatch{cOtherCrop} & other\_crop \\
10 & winter\_triticale         & 1 & \swatch{cWheat}     & wheat (pragmatic) \\
11 & winter\_durum\_wheat      & 1 & \swatch{cWheat}     & wheat \\
12 & fruits\_vegetables\_flowers & 6 & \swatch{cOtherCrop} & other\_crop \\
13 & potatoes                  & 6 & \swatch{cOtherCrop} & other\_crop \\
14 & leguminous\_fodder        & 5 & \swatch{cGrassland} & grassland \\
15 & soybeans                  & 6 & \swatch{cOtherCrop} & other\_crop \\
16 & orchard                   & 6 & \swatch{cOtherCrop} & other\_crop \\
17 & mixed\_cereal             & 1 & \swatch{cWheat}     & wheat \\
18 & sorghum                   & 3 & \swatch{cCorn}      & corn \\
19 & void\_label               & 7 & \swatch{cVoid}      & void (ignore index) \\
\bottomrule
\end{tabular}
\end{table}

% -----------------------------------------------------------------------------
\begin{table}[h]
\centering
\caption{EuroCrops HCAT name $\to$ common 8-class (\texttt{hcat\_to\_common}).
Substring matching on the lower-cased \texttt{EC\_hcat\_n} attribute;
rules are evaluated in order and the first match wins. Polygons whose
\texttt{EC\_hcat\_n} matches no rule fall through to \texttt{other\_crop}.}
\label{tab:hcat-to-common}
\small
\setlength{\tabcolsep}{6pt}
\begin{tabular}{C C l L}
\toprule
order & color & common class & substring(s) tested in \texttt{EC\_hcat\_n} \\
\midrule
1 & \swatch{cWheat}     & wheat       & \texttt{wheat}, \texttt{triticale} \\
2 & \swatch{cBarley}    & barley      & \texttt{barley} \\
3 & \swatch{cCorn}      & corn        & \texttt{maize}, \texttt{corn} \\
4 & \swatch{cRapeseed}  & rapeseed    & \texttt{rape} \\
5 & \swatch{cGrassland} & grassland   & \texttt{grass}, \texttt{meadow}, \texttt{pasture}, \texttt{lawn} \\
6 & \swatch{cNonCrop}   & non\_crop   & \texttt{fallow}, \texttt{urban}, \texttt{water}, \texttt{forest}, \texttt{road}, \texttt{wood}, \texttt{tree} \\
7 & \swatch{cOtherCrop} & other\_crop & (default for any other agricultural label) \\
\bottomrule
\end{tabular}
\end{table}

\paragraph{Notes.}
\begin{itemize}
  \item PASTIS class 10 (\texttt{winter\_triticale}) is pragmatically folded
        into \texttt{wheat} rather than \texttt{other\_crop}; triticale is a
        wheat--rye hybrid and the trained model's recall on the dedicated
        triticale class was poor.
  \item PASTIS class 14 (\texttt{leguminous\_fodder}) is mapped to
        \texttt{grassland} on the prediction side because in NRW such parcels
        are visually and phenologically similar to permanent grassland.
  \item Class 7 (\texttt{void}) is used as the loss/ignore index on the
        EuroCrops side, ensuring that pixels without a labelled parcel do not
        contribute to overall pixel accuracy, IoU, precision, or recall.
\end{itemize}

\end{document}
